\scanpage{183}
% manuscript position: top
\mspage{99}
% manuscript position: 1–2
ordonné grand devant $\pi$. Mais ceci va résulter du résultat plus précis :

% manuscript position: 3–7
\uline{Corollaire} 4.4.15. Soit $C$ une petite catégorie. Alors \inserted{la sous-catégorie strictement pleine} $\operatorname{Ind}'_\pi(C)$ \deleted{\unsure{est une sous-}} de $\widehat C$ est stable \inserted{dans $\widehat C$} par $\varinjlim$ grandes devant $\pi$, \deleted{\unsure{en fait}} \deleted{et} \inserted{en d'autres termes} ces $\varinjlim$ \deleted{\unsure{nt}}\unsure{existent|existe} \uline{dans $\operatorname{Ind}'_\pi(C)$}, et le foncteur d'inclusion y \unsure{commute}.
% manuscript position: left/4–7
% slanted
\margin{\unsure{Comparer|Comp.} \unsure{à|avec} 4.7.3.1 \unsure{qui pourrait|qu'il y aurait} \unsure{pu|se} s'expliciter \unsure{ici}.}

% manuscript position: 8–13
\uline{Dém} Soit $I$ ens. ordonné, grand devant $\pi$, et \deleted{$\unsure{(F_i)}$} $i\longmapsto F_i$ un foncteur $I\to\operatorname{Ind}'_\pi(C)$, $F=\varinjlim_I F_i$ dans $\widehat C$, prouvons que $F\in\operatorname{Ob}\operatorname{Ind}'_\pi(C)$, i.e. (4.4.3) que $C_{/F}$ est grd devant $\pi$. J'utilise le critère b) de 4.4.4, il faut vérifier $\mathrm{PF}_{\pi}1$, $\mathrm{PF}_{\pi}2$. Soit
% manuscript position: 14
\[
x_\alpha\xrightarrow{u_\alpha}F=\varinjlim F_i
\]
% manuscript position: 15–23
des objets de $C_{/F}$ indexés par $A$ avec $\operatorname{card} A\leq\pi$. Chaque $u_\alpha$ provient de $x_\alpha\to F_{i_\alpha}$, et en prenant un majorant commun des $i_\alpha$, \inserted{($I$ grd devant $\pi$)} on trouve un $i\in I$ et des $v_\alpha:x_\alpha\to F_i$. Comme $F_i$ est \deleted{\unsure{dans}} $\operatorname{Ind}'_\pi(C)$ i.e. $C_{/F_i}$ grd devant $\pi$, considérant les $v_\alpha$ comme des objets de $C_{/F_i}$, ils sont majorés par un objet $x\longrightarrow F_i$ de $F_i$. On a donc
% manuscript position: 24–25
\diagram{\begin{tikzcd}[ampersand replacement=\&]
x_\alpha \arrow[r] \arrow[d] \& F_{i_\alpha} \arrow[d] \arrow[dr,dashed] \& {}\\
x \arrow[r] \& F_i \arrow[r] \& F\text{,}
\end{tikzcd}}
% manuscript position: 26
qui montre que les $x_\alpha\longrightarrow F$ sont majorés par $x\to F$,
% manuscript position: foot
% archive sheet number: 182
